\documentclass[12pt]{article}

\usepackage[portuguese]{babel}
\usepackage{float}
\usepackage{listings}
\usepackage{indentfirst}
\usepackage{amsfonts, amssymb, amsmath}
\usepackage{graphicx}
\usepackage{enumitem}
\usepackage[colorlinks=true, allcolors=blue]{hyperref}
\usepackage[letterpaper,
            top=1.5cm,
            bottom=1.5cm,
            left=2.75cm,
            right=2.75cm,
            marginparwidth=1.75cm]{geometry}

\title{Organização e Recuperação de Dados}
\author{Prof. Jonathan}
\date{\today}

\begin{document}

\maketitle

\tableofcontents
\break

\section{Memória Secundária}
\begin{itemize}
    \item Por que discos e fitas são lentos?
    \begin{itemize}
        \item Por que eles são dispositivos eletromagnéticos que envolvem 
        partes mecânicas, enquanto outros tipos de memoria são eletrônicas.
    \end{itemize}
    \item SSDs não possuem partes mecânicas, por isso são dispositivos 
    secundários mais rápidos
    \begin{itemize}
        \item Mas ainda são milhares de vezes mais lentos do que a memória RAM
    \end{itemize}
    \item Fitas são dispositivos seriais, enquanto isso, HDs e SSDs são 
    dispositivos de acesso aleatorio
\end{itemize}

\subsection{Discos Magnéticos}
    \begin{itemize}
        \item Composição:
        \begin{itemize}
            \item Um ou mais pratos circulares sobrepostos atravessados por um 
            eixo que os rotaciona
            \item Os pratos são construidos com materiais não magnéticos e 
            revestidos por uma cobertura magnetica extremamente fina
            \item Cabeças de leitura/escrita posicionadas sovre as superficies 
            do prato sçao responsaveis pela leitura e gravação dos dados
            \item \textbf{SEEKING} é um termo importante
        \end{itemize}
        \item Funcionamento
        \begin{itemize}
            \item Durante uma leitura/gravação a cabeça permanece parada 
            enquanto o prato gira sob ela
            \begin{itemize}
                \item Esta diposição da origem a organizalçao dos dados no 
                prato como um conjunto de aneis concentricos chamados de trilhas
            \end{itemize}
            \item Pulsos magneticos enviados a uma bobina na cabeça de L/E 
            produzem campos magneticos que alteram a polaridade de pequenas 
            regioes de uma trilha do prato
            \begin{itemize}
                \item A diferença no sentido da polaridade entre cada região 
                codifica a informação binaria (0s e 1s) da trila 
            \end{itemize}
            \item Durante a leitura, quando as regioes magnetiadas passam sob 
            a cabeça, mudanças de polaridade induzem pulsos elétricos na bobina 
            da cabeca são convertidos novamente para a informação binária 
            equivalente
        \end{itemize}
        \item Organização dos discos
        \begin{itemize}
            \item Os dados são armazenados em trilhas sucessivas na superfície 
            de um prato
            \item Cada trilha é dividida em setores
            \begin{itemize}
                \item Um setor é a menor parte endereçável do disco (512B/4KB)
            \end{itemize}
            \item Quando é realizada uma chamada READ() por um bit específico
            \begin{itemize}
                \item O disco localiza a superficie, a trilha e o setor corretos
                \item Os dados de um setor inteiros são copiados para um buffer
                \item A partir do buffer, se encontra o byte especifico que foi 
                requisitado
            \end{itemize}
            \item Normalmente os discos vêm setorizados de fabrica (formatação 
            física)
            \item OBS.: É importante notar que, lógicamente, o disco é visto 
            como um vetor continuo de setores, i.e, a organização fisica é 
            transparente para o sistema
            \item Lacunas:
            \begin{itemize}
                \item Lacunas entre setores: Evitam que seja imposta uma precisão 
                absurda para a cabeça de L/E
                \item Lacunas entre trilhas: Evitam, ou pelo menos minimizam, 
                erros devido ao desalinhamento da cabeça ou simplesmente à 
                interferencia de campos magnéticos 
            \end{itemize}
            \item Em um disco formado por diversos pratos, o conjunto de 
            trilhas na mesma direção (sobrepostas em pratos diferentes) forma 
            um cilindro
            \begin{itemize}
                \item Um prato possui o mesmo número de trilhas na parte de 
                cima, e na parte de baixo. (Na hora de contar as trilhas de um 
                prato, se contar só pela parte de cima, lembrar de multiplicar 
                por 2)
            \end{itemize}
            \item Toda a informação adicional em um cilindro pode ser acessada 
        \end{itemize}
        \item Estimativa de capacidade
        \begin{itemize}
            \item A quantidade de dados que pode ser armazenada em uma trilha 
            depende do quão densamente os bits podem ser armazenados
            \item A densidade depende da qualidade da midia e da precisão das 
            cabeças de leitura/escrita
            \begin{itemize}
                \item Ex.: Um disco rigido de 1TB Seagate 7200RPM ($\pm$ R\$300,00)
                \begin{itemize}
                    \item 16383 trilhas/superfície e 63 setores (de 4k) por 
                    trilha (256KB/trilha)
                \end{itemize}
            \end{itemize}
            \item Sendo que: 
            \begin{itemize}
                \item Números de Cilindros = Número de trilhas/superfície
                \item Número de Trilhas do Cilindro = 2x o número de pratos 
            \end{itemize}
            \item Temos que:
            \begin{itemize}
                \item Capacidade da trilha = 
                \[
                \text{Numero de bytes por setor} \times \text{Números de setor 
                por trilha} 
                \]
                \item Capacidade do cilindro = 
                \[
                \text{Capacidade da trilha} \times \text{Números de trilhas do 
                cilindro} 
                \]
                \item Capacidade do disco = 
                \[
                \text{Capacidade do cilindro} \times \text{Número de cilindros} 
                \]
            \end{itemize}
            \item Exemplo
            \begin{itemize}
                \item Propriedades do disco:
                \begin{itemize}
                    \item 512 byts/setor
                    \item 63 setores/trilha
                    \item 16 trilhas/cilindro
                    \item 4080 cilindros
                \end{itemize}
                \item Quantos cilindros são necessários para armazenar um 
                arquivo de 50000 registros de 256 bytes cada
                \begin{itemize}[label=R:]
                    \item \dots

                \end{itemize}
            \end{itemize}
        \end{itemize}
        \item Clusters
        \begin{itemize}
            \item Organização lógica que visa aumentar o desempenho e mantida 
            pelo gerenciador de arquivos, presente no S.O
            \begin{itemize}
                \item Quando um arquivo é acessado por uma aplicalçao é o 
                gerenciador de arquivos quem associa o arquivo lógico às suas 
                posições fisicas
                \item O arquivo é visto como uma série de clusters
            \end{itemize}
            \item Para fazer esse mapeamento, o gerenciador de arquivos
            utiliza uma tabela de alocação de arquivos (chamada de File 
            Allocation Table (FAT) em alguns S.Os.)
            \item Em vez da tabela de alocação endereçar setores, endereça 
            clusters
            \begin{itemize}
                \item Um cluster é um conjunto de um ou mais setores contíguos 
                do disco
                \begin{itemize}
                    \item Todos os setores de um cluster podem ser lidos sem 
                    seeks adicionais e sem a necessidade de consultas adicionais 
                    à tabela de alocação
                \end{itemize}
                \item A tabela de alocação contem uma lista de todos os 
                clusters de um arquivo, ordenados de acordo com a ordem lógica 
                dos setores que eles contém
            \end{itemize}
            \item Cada cluster do disco é usado para um unico arquivo, ou seja, 
            em um mesmo cluster não haverá informações sobre mais de um arquivo. 
        \end{itemize}
        \item Custos de acesso a disco
        \begin{itemize}
            \item Em termos de tempo, o custo de um acesso ao disco é a soma 
            de três tempos:
            \begin{itemize}
                \item Posicionamento (Seeking) -- Tempo necessário para mover a 
                cabeça de L/E até a trilha correta (Depende da distância a ser 
                percorrida)
                \item Latência (Rotational Delay ou Latency) -- Tempo necessário 
                para a cabeça de L/E se posicionar no setor correto.
                \item Transferência (Transfer Time) -- Tempo necessário para um 
                byte ser lido na superfície do disco e transferido para o buffer 
                interno da controladora
            \end{itemize}
        \end{itemize}
    \end{itemize}

    \subsection{Custo de Acesso ao Disco}
    \begin{itemize}
        \item Seek Time -- Posicionamento
        \begin{itemize}
            \item É impossivel saber exatamente quantas trilhas serão 
            percorridas durante as buscas
            \item O que se faz é estimar um tempo médio, assumindo que as 
            posições iniciais e finais para cada acesso são aleatórias
            \item Por meio de estudos empiricos, estumou-se que uma busca 
            percorre, em média, 1/3 do total de trilhas, sendo que o tempo 
            gasto para percorrer esse numero de trilhas tem sido usado 
            pelos fabricantes como indicador do tempo medio de busca (seek 
            time).
            \item O que ocorre quando um arquivo está armazenado em 
            cilindros consecutivos e é acessado sequencialmente?
            \begin{itemize}
                \item O seek time é reduzido! (Melhor situação)
            \end{itemize}
            \item O que ocorre quando dois arquivos, localizados em
            extremos opostos do disco (um no cilindro mais externo e
            outro no mais interno), são acessados alternadamente?
            \begin{itemize}
                \item O seek time é alto! (Pior situação)
            \end{itemize}
            \item O tempo de seek tende a ser mais caro em sistemas
            multiusuários em que vários processos concorrem pelo uso do 
            disco
        \end{itemize}
        \item Latência
        \begin{itemize}
            \item Tempo gasto para a cabeça de L/E encontre o setor desejado
            \begin{itemize}
                \item Estima-se que o tempo medio seja a metade do tempo 
                de uma rotação (na pratica, esse tempo costima ser menor)
                \item Tambem chamado de Avarage Latency
                \item No inicio dos anos 90, com discos de 3600rpm, a 
                metade do tempo de rotação era de 8,3ms 
                \item Atualmente os discos de 7200rpm tem latencia media 
                de 4,16ms
            \end{itemize}
            \item Em uma leitura/escrita sequencial envolvendo trilhas de 
            um mesmo cilindro, não há latencia na alternânacia entre as 
            trilhas
            \begin{itemize}
                \item Existe apenas um tempo (muito pequeno) de comutação 
                entre as cabe-ças de leitura/escrita
            \end{itemize}
        \end{itemize}
        \item Tempo de Transferência
        \begin{itemize}
            \item Uma vez que a cabeça de L/E está sob o setor desejado, 
            ele pode ser transferido
            \item O tempo de transferencia é dado por:
            \begin{itemize}
                \item O tempo para transferir uma trilha inteira normalmente
                é o tempo de rotação
                \item ($\mathrm{n^o}$ bytes transferidos/$\mathrm{n^o}$ 
                bytes na trilha) $\times$ tempo de rotação
                \item Exemplo: 
                \begin{itemize}
                    \item Tempo de transferencia de 1KB em disco com 
                    32 setores de 512 bytes por trilha (16.384b) e tempo de 
                    rotação de 8,2ms \[(1024/16384) \times 8,2 = 0,51ms\]
                    \item Neste exemplo o tempo está expresso em 
                    bytes/ms. Os fabricantes costumam utilizar MB/s
                \end{itemize}
            \end{itemize}
        \end{itemize}
        \item O modo de acesso  ao arquivo pode afetar drasticamente os 
        custos de tempo de acesso
        \begin{itemize}
            \item Dois modos de acesso:
            \begin{itemize}
                \item Sequencial: o máximo do arquivo é processo em cada 
                acesso
                \item Aleatório: apenas um registro é processado por vez
            \end{itemize}
            \item Exemplo:
            \begin{itemize}
                \item Determinar o tempo necessário para ler um arquivo com 
                40000 registros de 256 bytes para cada um
                \begin{itemize}
                    \item Vamos calcular o tempo de leitura do arquivo com 
                    acesso sequencial e aleatório e comparar resultados
                \end{itemize}
                \item Vamos considerar as seguintes especificações do disco:
                \begin{table}[H]
                    \centering
                    \begin{tabular}{|c|c|}
                        \hline
                        Tempo médio de seek & 13 ms \\ \hline
                        Tempo de latência & 8.3 ms \\ \hline
                        Tempo de transferencia & 16.7 ms ou 1229 bytes/ms \\
                        \hline
                        Bytes por setor & 512 \\ \hline
                        Setores por trilha & 100 \\ \hline
                        Trilhas por cilindro & 12 \\ \hline
                        Trilhas por superficie & 1748 \\ \hline
                        Tamanho do cluster & 10 setores (5120 bytes) \\ \hline
                    \end{tabular}
                \end{table}
                \item Tamanho do arquivo $Rightarrow$ 40.000 registros de 256 
                bytes
                \begin{itemize}
                    \item Se cada cluster tem 5120 bytes (10 setores), podemos 
                    armazenar 20 registros por cluster (5120/256)
                    \item Será necessário um total de 2000 clusters para 
                    armazenar todos os registros (40000/20)
                    \item Como cada trilja possui 10 clusters (100 setores), 
                    serão necessárias 200 trilhas
                \end{itemize}
                \item Vamos assumir o pior caso, no qual as 200 trilhas que 
                armazenam o arquivo estão espalhadas aleatoriamente no disco
                \begin{itemize}
                    \item Situação extrema, mas que pode ocorrer em discos no 
                    limite da capacidade, especialmente com muitos arquivos 
                    pequenos
                \end{itemize}
                \item Custo com acesso sequencial
                \begin{itemize}
                    \item Para cada trilha em que são lidos setores consecutivos, 
                    o processo de leitura envolve os seguintes custos: Tempo 
                    médio de seek, tempo de latência e tempo de transferencia de 
                    uma trilha.
                    \item Para 200 trilhas: 
                    \[(13 + 8,3 + 16,7) \times 200 = 7600ms = 7,6s\]
                \end{itemize}
                \item Custo com acesso aleatório
                \begin{itemize}
                    \item Para cada trilha em que são lidos setores consecutivos, 
                    o processo de leitura envolve os seguintes custos: Tempo 
                    médio de seek, tempo de latência e tempo de transferencia de 
                    um cluster.
                \end{itemize}
            \end{itemize}
        \end{itemize}
    \end{itemize}
    
    \subsection{SSD}
    \begin{itemize}
        \item Solid-state drives(SSD) são compostos por componentes eletrônicos
        \begin{itemize}
            \item O armazenamento geralmente é feito em módulos de memória flash.
        \end{itemize}
        \item Devido a essa caracteristica, o custo de acesso a qualquer 
        posição do SSD será o mesmo
        \begin{itemize}
            \item Não há tempo de seek, nem latencia rotacional
            \item Mas existe uma latência (tempo para se iniciar uma operação) 
            e um tempo de transferencia (throughput)
            \begin{itemize}
                \item Esses tempos dependem do tipo de memória usada, da 
                controladora e da interface do SSD
            \end{itemize}
        \end{itemize}
        \item A menor unidade endereçável de um SSD é uma página.
        \begin{itemize}
            \item O tamanho de uma página muda de um SSD para outro, variando 
            entre 2KB e 16KB.
            \item Assim como acontece com os setores de um disco, a pagina é a 
            menor unidade de leitura e escrita.
        \end{itemize}
        \item As paginas são agrupadas em blocos.
        \begin{itemize}
            \item O tamanho dos blocos tambem varia de um SSD para outro, 
            entre 256KB e 4MB 
            \begin{itemize}
                \item Por exemplo, o Samsung SSD 840 EVO tem blocos de 2048KB 
                (256 páginas de 8KB cada)
            \end{itemize}
        \end{itemize}
        \item Diferentemente dos setores de um HD, as paginas de um SSD não 
        podem ser sobrescritas.
        \begin{itemize}
            \item Quando uma alteração é necessaria, a pagina é copiada para um 
            buffer, alterada a escrita em uma nova página livre.
            \item A página antiga é marcada como "velha" e ficará assim até o 
            seu bloco ser apagado e a página fique livre novamente.
        \end{itemize}
        \item Essa operação faz com que as escritas sejam mais lentas que as 
        leituras.
        \item Alem disso, apenas blocos podem ser apagados.
        \begin{itemize}
            \item Para que um bloco com alguma página "velha" seja apagado, o 
            coletor de lixo copiará as paginas ativas para um novo bloco, 
            deixando o bloco pronto para ser apagado por completo.
            \item Uma vez apagado, as páginas do bloco ficam livres novamente.
        \end{itemize} 
        \item Além de não poderem ser sobrescritas, memórias flash também têm 
        um numero máximo de escritas possiveis.
        \begin{itemize}
            \item Elas desgastam conforme são escritas e apagadas
            \item Por conta disso, a controladora distribuirá o uso dos blocos 
            de memoria para que o desgaste não seja maior em uma região do que 
            em outra.
        \end{itemize}
        \item Todos esses fatores fazem com que as controladoras sejam complexas 
        e adicionem custo nas operações em termos de tempo, especialmente nas 
        escritas.
        \item Ainda assim, os SSDs são centenas de vezes mais rapidos do que os 
        HDs, especialmente quando se trata de acesso aleatório.
    \end{itemize}
\vspace{24pt}
Independente do tipo de dispositivo usado para como memória secundária, ele será
o gargalo do sistema.
\begin{itemize}
    \item Gargalo $\rightarrow$ Quando um dispositivo mais lento afeta afeta o 
    desempenho de outros mais rapidos, se tornando um fator limitante do sistema.
\end{itemize}
Por isso é importante que acessos à memória secundária sejam feitos de forma 
eficiente.

\section{Operações em Arquivos}
    \subsection{Arquivos}
    \begin{itemize}
        \item Um arquivo é uma estutura criada em uma unidade de armazenamento 
        secundário
        \begin{itemize}
            \item Permite o armazenamento permanecente dos dados, ao contratio 
            das variáveis, que são armazenadas em RAM
        \end{itemize}
        \item Do ponto de vista físico, um arquivo corresponde a uma sequencia 
        de bytes armazenados em um dispositivo secundário
        \begin{itemize}
            \item Mas os aspectos fisicos dos arquivos são transparentes para 
            as aplicações
        \end{itemize}
    \end{itemize}
        \subsubsection{Diferentes Visões do Arquivo}
        \begin{itemize}
            \item Arquivo físico e arquivo lógico
            \begin{itemize}
                \item Arquivo físico: é o arquivo do ponto de vista do 
                armazenamento
                \begin{itemize}
                    \item Corresponde a uma sequencia particularde bytes 
                    armazenados no dispositivo
                \end{itemize}
                \item Arquivo lógico: é o arquivo do ponto de vista da aplicação 
                que o acessa
                \begin{itemize}
                    \item Cada arquivo manipulado pela aplicação é tratado como 
                    um canal de comunicação estabelecido entre a aplicação e o 
                    arquivo físico
                \end{itemize}
            \end{itemize}
            \item O sistema operacional é quem faz a interface entre o arquivo 
            lógico e o arquivo físico.
            \item Apesar de um dispositivo secundário poder conter milhares de 
            arquivos físicos, as aplicações, em geral, só podem abrir um número 
            pré-definido de arquivos lógicos de forma simultânea
            \item Um arquivo lógico pode estar associado a um arquivo físico ou 
            a outros dispositivos de E/S, como o teclado (stdin) e o vídio 
            (stdout e stderr)
            \item Para os sistemas operacionais, um dispositivo de E/S é 
            visto da mesma forma que um arquivo!
        \end{itemize}
        \subsubsection{Conectando Arquivo Lógico ao Físico}
        \begin{itemize}
            \item A aplicação solicita que um nome lógico seja associado a um 
            arquivo físico específico e o S.O estabelece a associação


        \end{itemize}
    \subsection{Arquivo Lógico}
    \begin{itemize}
        \item Do ponto de vista da aplicação, um arquivo é uma entidade 
        lógica vista como uma sequência de bytes
        \begin{itemize}
            \item Podemos enviar bytes para o arquivo indefinidamente
            \item De fomra análoga, podemos ler bytes do arquivo enquanto 
            houverem bytes a serem lidos
        \end{itemize}
        \item Para todo arquivo lógico, é mantido um ponteiro de L/E que 
        indica a posição do próximo byte a ser lido/escrito
        \begin{itemize}
            \item Esse ponteiro se move automáticamente conforme operações 
            de leitura e escrita são realizadas e também pode ser movido 
            por comandos de reposicionamento.
        \end{itemize}
        \item Quando um arquivo é aberto, o ponteiro de L/E é posicionado 
        no início (byte 0)

        \item Se for feita a leitura de 9 bytes, o ponteiro de L/E se move 
        automaticamente para a proxima posição de L/E (byte 9) 

        \item Se for feita uma escita de 8 bytes na sequencia, o ponteiro 
        de L/E se moverá para a próxima posição de L/E (byte 17)

        \item O arquivo lógico é identificado pelo S.O por um número 
        inteiro chamado de descritor
        \item É comum que as linguagens encapsulem o descritor do arquivo 
        em estruturas mais complexas, que armazenam outras informações úteis 
        para a manipulação do arquivo
        \begin{itemize}
            \item O tipo dessa estrutura descritora é dependente da linguagem
            \item Em Python, essas informações ficam em um objeto de arquivo
        \end{itemize}
    \end{itemize}
    \subsection{Tipos de Arquivos}
    \begin{itemize}
        \item Basicamente são dois tipos de arquivos
        \begin{itemize}
            \item Texto
            \begin{itemize}
                \item Os caracteres são armazenados sequencialmente
                \item É possivel determinar o primeiro, segundo, terceiro, 
                $n$ caractéres que com-põem o arquivo
                \item Nesse tipode arquivo, a sequencia 123 correspode à string 
                "123"
            \end{itemize}
            \item Binário
            \begin{itemize}
                \item Formado por uma sequencia de bytes sem correspondencia 
                direta com um tipo de dado
                \item Cabe ao programador fazer essa correspondencia quando lê 
                e escreve nesses programas
                \item Nesse tipo de arquivo, o valor 123 será gravado como um 
                numero binario com uma quantidade pré-definida de bytes
            \end{itemize}
            \item Exemplo:
            \begin{itemize}
                \item As figuras mostram um arquivo binário e um arquivo texto 
                em um editor hexadecimal
            \end{itemize}
        \end{itemize}
    \end{itemize}
    \subsection{Operações em arquivos}
    \begin{itemize}
        \item Basicamente as operações em arquivos são as seguintes:
        \begin{enumerate}[label=\arabic*.]
            \item Criar um arquivo
            \item Abrir um arquivo
            \item Ler de um arquivo
            \item Escrever em um arquivo
            \item Movimentar o ponteiro de L/e
            \item Fechar um arquivo
        \end{enumerate}
        \item Veremos como essas operações são feitas em Python
        \item As funções da linguagem Python para manipular arquivos são 
        implementadas de duas formas
        \begin{itemize}
            \item Como função nativa
            \item Como métodos de objetos de arquivo
        \end{itemize}
        \item A função open é nativa e utilizada para abrir/criar um arquivo 
        retornando um objeto de arquivo
        \begin{itemize}
            \item A classe desse objeto varia de acordo com o tipo de arquivo a 
            ser manipulado (leitura/escrita textual, leitura/escrita binária, 
            etc.)
            \item A partir do objeto de arquivo é possivel acessar métodos para 
            leitura, escrita, posicionamento, etc.
        \end{itemize}
        \item A associaão entre arquivo lógico e íssico é feita no momento da 
        abertura/criação do arquivo.
        \begin{itemize}
            \item O modo como associação se dará depende do parametro utilizado 
            para especificar o modo de abertura
            \item O modo de abertura também especifica se um novo arquivo será 
            criado ou se um arquivo existente será aberto
        \end{itemize}
    \end{itemize}
        \subsubsection{Criar e abrir arquivos}
        \begin{itemize}
            \item \verb|def open(filename:str, mode:str) -> file object|
            \begin{itemize}
                \item filename é uma string contendo o nome do arquivo físico
                \item mode também é uma string e define o modo de abertura
                \item \textbf{Caso a operação falhe, um erro ocorrerá e abortará 
                a execução!}
                \item Exemplo: \verb|arq = open('meuarq.txt','w')|
                \item Flags/modos da função \verb|open()|
                \begin{table}[H]
                    \centering
                    \begin{tabular}{|c| p{10cm}|}
                        \hline
                        Caractere & Descrição do comportamento \\ \hline
                        'r' & Abre para leitura. O arquivo deve existir. 
                        \textbf{É o modo padrão}. \\ \hline
                        'w' & Abre para escrita. Se o arquivo existir ele será 
                        truncando. Senão, será criado. \\ \hline
                        'x' & Abre um novo arquivo para escrita. O arquivo não 
                        deve existir. \\ \hline
                        'a' & Abre para escrita no final do arquivo. Operações 
                        de reposicionamento do ponteiro de L/E são ignoradas. 
                        Se não existir, o arquivo será criado. \\ \hline
                        't' & Abre em modo texto. Admite apenas dados do tipo 
                        \verb|str|. \textbf{É o modo padrão}. \\ \hline
                        'b' & Abre em modo binário. Admite apenas dados do tipo 
                        bytes/bytearray. \\ \hline
                        '+' & Altera o comportamento da flag inicial para 
                        permitir leitura e escrita. \\ \hline
                    \end{tabular}
                \end{table}
                \begin{itemize}
                    \item As flags 't','b' e '+' só podem ser utilizadas em 
                    conjunto de outras!
                    \item Por padrão os arquivos são abertos em modo texto 
                    (flag 't')
                \end{itemize}
            \end{itemize}
            \item Como tratar um erro na abertura do arquivo?
            \begin{itemize}
                \item Em Python, erros de execução são chamados de exceções
                \item As exceções podem ser capturadas e tratadas utilizando a 
                instrução try-except
                \begin{verbatim}
    try:
        ''' código a ser executado '''
    except:
        ''' código a ser executado caso uma exceção ocorra '''
                \end{verbatim}
                \item Exemplo:
                \begin{verbatim}
    try:
        arq = open('meuArq.txt','r')
        print('abertura OK')
    except:
        print('Erro na abertura')
        exit()
                \end{verbatim}
            \end{itemize}
        \end{itemize}
        \subsubsection{Fechamento de arquivos}
        \begin{itemize}
            \item Uma vez que o arquivo foi aberto, seu fechamento se dá pela 
            chamada do método \verb|close()|
            \begin{itemize}
                \item Todos os objetos de arquivo possuem esse método
                \item Encerra a associação entre o arquivo lógico e o físico
                \item O descritor fica disponível para uso por outro arquivo
                \item Garante que todas as informações do arquivo foram 
                atualizadas (o fechamento garante que o conteudo bufferizado 
                foi enviado para o dispositivo físico) 
                \item O S.O normalmente fecha qualquer arquivo que ficou aberto 
                quando um programa termina corretamente, mas o ideal é fechar 
                arquivos que não estão mais sendo utilizados
                \begin{itemize}
                    \item Previne a perda de dados bufferizados no caso de erro
                    \item Libera a estrutura descritora para novos arquivos
                \end{itemize}
            \end{itemize}
            \item Exemplo:
            \begin{itemize}
                \item \begin{verbatim}
arq = open('meuarq.txt', 'r')
''' comandos '''
arq.close()\end{verbatim}
            \end{itemize}
            \item Quando o arquivo é aberto em um comando \verb|with|, não é 
            necessário o fechamento explícito
            \begin{itemize}
                \item \begin{verbatim}
with open('meuarq.txt', 'r') as arq
    ''' comandos '''\end{verbatim}
                \begin{itemize}
                    \item Nesse caso, o objeto de arquivo arq só existirá no 
                    escopo do comando with
                \end{itemize}
            \end{itemize}
        \end{itemize}
        \subsection{Leitura e escrita em arquivos}
        \begin{itemize}
            \item Em Python, as funções de leitura e escrita são métodos do 
            objeto de arquivo
            \begin{itemize}
                \item Portanto, dependem do tipo do objeto
            \end{itemize}
            \item Arquivos em modo texto são representados por objetos do tipo 
            \textbf{TextIOWrapper}
            \begin{itemize}
                \item Só admitem leitura e escrita de strings, i.e., objetos 
                do tipo str
            \end{itemize}
            \item Arquivos em modo binário são representados por objetos 
            \textbf{BufferedReader} (leitura), \textbf{BufferedWriter} (escrita) 
            e \textbf{BufferedRandom} (leitura e escrita)
            \begin{itemize}
                \item Só admitem leitura e escrita de objetos do tipo bytes ou 
                bytearray
                \begin{itemize}
                    \item bytes são imutaveis e bytearray são mutaveis -- ambos 
                    representam sequencias de bytes
                \end{itemize}
            \end{itemize} 
        \end{itemize}
            \subsubsection{Leitura em modo texto}
            \begin{itemize}
                \item \verb|def read(size=-1) -> str|
                \begin{itemize}
                    \item Lê e retorna no maximo size caracteres como uma unica 
                    string
                    \item Se size for negativo ou None, lê até o fim do arquivo 
                    (EOF), i.e., retorna todo o conteudo do arquivo como um 
                    unica string
                \end{itemize}
                \item \verb|def readline(size=-1) -> str|
                \begin{itemize}
                    \item Lê até uma quebra de linha ou EOF e retorna o conteudo 
                    como uma unica string
                    \item Se o ponteiro L/E já estiver EOF, retorna uma string 
                    vazia
                    \item Se size for especificado, no maximo size caracteres 
                    serão lidos
                \end{itemize}
                \item \verb|def readlines(hint=-1) -> list[str]|
                \begin{itemize}
                    \item Retorna as linhas do arquivo como uma lista de strings
                    \item Se hint for especificado, linhas serão lidas até o 
                    maximo de hint caracteres no total
                \end{itemize}
            \end{itemize}
            \subsubsection{Escrita em modo texto}
            \begin{itemize}
                \item \verb|def write(s) -> int|
                \begin{itemize}
                    \item Escreve a string s no arquivo
                    \item Retorna o numero de caracteres escritos
                \end{itemize}
                \item \verb|writelines(lines) -> None|
                \begin{itemize}
                    \item Escreve uma lista de strings
                    \item Não são adicionados separadores entre as strings, 
                    então é comum que cada string contenha um caractere 
                    separador (p.e., '\verb|\n|') no final.
                \end{itemize}
            \end{itemize}
            \subsubsection{Leitura em modo binário}
            \begin{itemize}
                \item \verb|def read(size=-1) -> bytes|
                \begin{itemize}
                    \item Lê e retorna até size bytes
                    \item Se size for negativo ou None, o conteudo do arquivo 
                    será lido até EOF e retornado 
                    \item Um objeto bytes vazio(b'') é retornado se o ponteiro 
                    estiver em EOF
                \end{itemize}
                \item \verb|def readinto(b) -> int|
                \begin{itemize}
                    \item Lê bytes para uma variavel b, pré-alocada como um 
                    bytearray, e retorna o numero de bytes lidos
                \end{itemize}
                \item \verb|def readlines(hint=-1) -> list[bytes]|
                \begin{itemize}
                    \item Retorna o conteudo do arquivo como uma lista de linhas
                    \item Considera o valor b'\verb|\n|' como delimitador de 
                    linha
                    \item Se hint for especificado, linhas serão lidas até o 
                    maximo de hint bytes no total
                \end{itemize}
            \end{itemize}
            \subsubsection{Escrita em modo binário}
            \begin{itemize}
                \item \verb|def write(b) -> int|
                \begin{itemize}
                    \item Escreve o objeto b do tipo bytes ou bytearray no 
                    arquivo
                    \item Retorna o numero de bytes escritos (sempre igual ao 
                    comprimento de b em bytes, pois uma falha na escrita irá 
                    gerar uma exceção)
                \end{itemize}
                \item \verb|def writelines(lines) -> None|
                \begin{itemize}
                    \item Escreve uma lista de linhas no arquivo
                    \item Não são adicionados separadores entre as linhas, então 
                    é comum que cada linha seja adicionada de um caractere 
                    separador (p.e., b'\verb|\n|')
                \end{itemize}
                \item Detecção de fim de arquivo
                \begin{itemize}
                    \item Durante a leitura de um arquivo, o S.O monitora a 
                    posição do ponteiro de L/E do arquivo 
                \end{itemize}
            \end{itemize}
            
\section{Acesso Sequencial \dots}

\section{Busca Binária}
\begin{itemize}
    \item \dots
    \item \dots
\end{itemize}
    \subsection{Busca Binária em memória}
    \begin{itemize}
        \item Função que recebe um valoe  a ser buscado e uma lista ordenada 
        valoe
        \item Retorna \dots
    \end{itemize}
    \subsection{Busca binária em um arquivo}
    \begin{itemize}
        \item Pre requisitos para BB em arquivo
        \begin{itemize}
            \item Requer registros de tamanho fixo
            \item O arquivo deve estar ordenado em ordem crescente da chave de 
            busca
        \end{itemize}
        \item Complexidade
        \begin{itemize}
            \item Em um arquivo de n registros
            \begin{itemize}
                \item Busca binaria: O($\log_2n$), pior caso $\lfloor \log_2n \rfloor$
                \begin{itemize}
                    \item Em cada chute elimina-se metade das possibilidades
                \end{itemize}
            \end{itemize}
        \end{itemize}
        \item A busca binaria é mais eficiente ddo que a sequencial, porem o 
        arquivo deve estar ordenado por chave
        \begin{itemize}
            \item Custo de ordenar e manter o arquivo ordenado após novas 
            inserções
            \begin{itemize}
                \item A cada inserção deve-se reordenar o arquivo ou fazer a 
                inserção no arquivo em memoria e regravar o disco
            \end{itemize}
        \end{itemize}
        \item Um primeira solução $\rightarrow$ Ordenação interna
        \begin{itemize}
            \item Só é fctivel se o arquivo for pequeno e couber na memória
            \begin{enumerate}
                \item Ler todo o arquivo do disco para a RAM (acesso sequencial)
                \item Ordenar os registros em RAM usando um algoritmo de 
                ordenação
            \end{enumerate}
        \end{itemize}
    \end{itemize}
\end{document}